Dans un examen, une épreuve notée sur dix points est constituée de deux exercices : le premier est noté sur deux points, le deuxième sur huit points.


\medskip

\noindent {\Large \textbf{Partie 1}}

\medskip




Le premier exercice est constitué de deux questions Q1 et Q2.

Chaque question est notée sur un point. Une réponse correcte rapporte un point; une réponse incorrecte, incomplète ou une absence de réponse rapporte zéro point.

\begin{list}{\textbullet}{On considère que :}
\item Un candidat pris au hasard a une probabilité $0,8$ de répondre correctement à la question Q1.
\item Si le candidat répond correctement à Q1, il a une probabilité $0,6$ de répondre correctement à Q2; s'il ne répond pas correctement à Q1, il a une probabilité $0,1$ de répondre correctement à Q2.
\end{list}

On prend un candidat au hasard et on note :

\begin{list}{\textbullet}{}
		\item $A$ l'évènement : \og{}le candidat répond correctement à la question Q1 \fg{} ;
		\item $B$ l'évènement : \og{}le candidat répond correctement à la question Q2 \fg{}.
	\end{list}

On note $\overline{A}$ et $\overline{B}$ les évènements contraires de $A$ et de $B$.

	\begin{enumerate}
		\item On complète l'arbre pondéré ci-dessous.

\begin{center}
{\bigskip
\psset{levelsep=3cm,nodesepB=4pt, treesep=1cm}
\pstree[treemode=R,nodesepA=0pt]
{\TR{}}
{
\pstree[nodesepA=4pt]{\TR{$A$}\naput{\blue $0,8$}}
{
	\TR{$B$}\naput{\blue $0,6$}
			\TR{$\overline{B}$}\nbput{\blue $0,4$}
	}
\pstree[nodesepA=4pt]{\TR{$\overline{A}$}\nbput{\blue $0,2$}}
	{
	\TR{$B$}\naput{\blue $0,1$}
			\TR{$\overline{B}$}\nbput{\blue $0,9$}
	}
}
}
\bigskip
\end{center}
		\item La probabilité que le candidat réponde correctement aux deux questions Q1 et Q2 est: $P(A\cap B) = P(A)\times P_{A}(B)= 0,8\times 0,6 = 0,48$
		\item La probabilité que le candidat réponde correctement à la question Q2 est $p(B)$. D'après la formule des probabilités totales:

$P(B) = P(A\cap B) + P \left (\overline{A}\cap B\right )=0,48+0,2\times 0,1=0,5$		

	\end{enumerate}

\begin{list}{\textbullet}{On note :}
		\item $X_{1}$ la variable aléatoire qui, à un candidat, associe sa note à la question Q1 ;
		\item $X_{2}$ la variable aléatoire qui, à un candidat, associe sa note à la question Q2 ;
		\item $X$ la variable aléatoire qui, à un candidat, associe sa note à l'exercice, c'est-à-dire $X = X_{1} + X_{2}$.
\end{list}

\begin{enumerate}
\setcounter{enumi}{3}

\item La loi de probabilité de $X_1$ est:\hspace{2cm}
$\begin{array}{|c|c|c|}
\hline
x_i & 0 & 1\\
\hline
P(X_1=x_i) & 0,2 & 0,8\\
\hline
\end{array}$

Donc $E(X_1)= 0\times 0,2 + 1\times 0,8=0,8$.

La loi de probabilité de $X_2$ est:\hspace{2cm}
$\begin{array}{|c|c|c|}
\hline
x_i & 0 & 1\\
\hline
P(X_2=x_i) & 0,5 & 0,5\\
\hline
\end{array}$

Donc $E(X_2)= 0\times 0,5 + 1\times 0,5 =0,5$.

$E(X)=E(X_1+X_2)=E(X_1)+E(X_2)$ d'après la linéarité de l'espérance mathématique.
Donc $E(X)=0,8+0,5=1,3$. Un élève obtient, en moyenne, la note de $1,3$ à l'exercice 1.

\item On souhaite déterminer la variance de $X$.

	\begin{enumerate}
		\item 
\begin{list}{\textbullet}{}
\item $P(X=0)=P\left (\overline{A}\cap \overline{B}\right )=0,2\times 0,9 = 0,18$
\item $P(X=2)=P(A\cap B)=0,48$
\item $P(X=1)=1-\left (P(X=0)+P(X=2)\strut\right )=1-(0,18+0,48)=1-0,66=0,34$
\end{list}

On peut donc en déduire la loi de probabilité de $X$:

\begin{center}
\begin{tabularx}{0.6\linewidth}{|c| *3{>{\centering\arraybackslash}X|}}
\hline
$x_i$ & 0 & 1 & 2\\
\hline
$P(X=x_i)$ & $0,18$ & $0,34$ & $0,48$\\
\hline
\end{tabularx}
\end{center}
		\item 
D'après la formule de König: $V(X)=E(X^2)-[E(X)]^2$.

$E(X^2)=0^2\times 0,18+1^2\times 0,34+2^2\times 0,48 = 2,26$

Donc $V(X)=2,26-1,3^2=0,57$
		\item
$V(X_1)= E(X_1^2)-[E(X_1)]^2 = \left (0^1\times 0,2+1^2\times 0,8\strut\right ) - (0,8)^2
= 0,8-0,64 = 0,16$

$V(X_2)= E(X_2^2)-[E(X_2)]^2 = \left (0^1\times 0,5+1^2\times 0,5\strut\right ) - (0,5)^2
= 0,5-0,25 = 0,25$

$V(X_1)+V(X_2)=0,16 + 0,25=0,41$; or $V(X)=0,57$ donc $V(X_1)+V(X_2)\neq V(X)$.

Ce n'est pas surprenant car les variables aléatoires $X_1$ et $X_2$ ne sont pas indépendantes.
		
	\end{enumerate}
\end{enumerate}


\medskip

\noindent {\Large \textbf{Partie 2}}

\medskip


Le deuxième exercice est constitué de huit questions indépendantes.

Chaque question est notée sur un point. Une réponse correcte rapporte un point ; une réponse incorrecte et une absence de réponse rapporte zéro point.

Les huit questions sont de même difficulté : pour chacune des questions, un candidat a une probabilité $\dfrac{3}{4}$ de répondre correctement, indépendamment des autres questions.

On note $Y$ la variable aléatoire qui, à un candidat, associe sa note au deuxième exercice, c'est-à-dire le nombre de bonnes réponses.

\begin{enumerate}
\item
On est dans le cas d'une répétition de 8 épreuves identiques et indépendantes, la probabilité du succès lors d'une épreuve étant $\dfrac{3}{4}$. Donc la variable aléatoire $Y$ qui compte le nombre de succès suit la loi binomiale de paramètres $n=8$ et $p=\dfrac{3}{4}$.

\item 
Pour une variable aléatoire $Y$ suivant la loi binomiale $\mathcal{B}(n,p)$, on a:\\
$P(Y=k)=\ds\binom{n}{k} p^k \left (1-p\right )^{n-k}$, 
donc $P(Y=8)=\ds\binom{8}{8} \left (\dfrac{3}{4}\right )^{8}\left (1-\dfrac{3}{4}\right )^{8-8}
= \left (\dfrac{3}{4}\right )^8$.

\item L'espérance de $Y$ est $E(Y)= np = 8\times \dfrac{3}{4}= 6$.

La variance de $Y$ est $V(Y)= np(1 - p) = 8\times \dfrac{3}{4}\times \left (1 - \dfrac{3}{4}\right )=\dfrac{6}{4}= \dfrac{3}{2}$.
\end{enumerate}

\medskip

\noindent {\Large \textbf{Partie 3}}

\medskip



On suppose que les deux variables aléatoires $X$ et $Y$ sont indépendantes. On note la variable aléatoire qui, à un candidat, associe sa note totale à l'examen : $Z = X + Y$.

\begin{enumerate}
\item 
$Z=X+Y$ donc $E(Z)=E(X)+E(Y)= 1,3+6=7,3$

$Z=X+Y$ et $X$ et $Y$ sont indépendantes donc $V(Z)=V(X)+V(Y)=0,57+1,5=2,07$

\item Soit $n$ un nombre entier strictement positif.

Pour $i$ entier variant de 1 à $n$, on note $Z_{i}$ la variable aléatoire qui, à un échantillon de $n$ élèves, associe la note de l'élève numéro $i$ à l'examen.

On admet que les variables aléatoires $Z_{1}, Z_{2}, \ldots, Z_{n}$ sont identiques à $Z$ et indépendantes.

On note $M_{n}$ la variable aléatoire qui, à un échantillon de $n$ élèves, associe la moyenne de leurs $n$ notes, c'est-à-dire :
$M_{n} = \dfrac{Z_{1} + Z_{2} + \cdots + Z_{n}}{n}$

	\begin{enumerate}
		\item 
Pour tout $i$ entre 1 et $n$, on a $E(Z_i)= 7,3$.

$E(M_n)=\dfrac{E(Z_{1}) + E(Z_{2}) + \cdots + E(Z_{n})}{n}
=\dfrac{n\times 7,3}{n} = 7,3$
		

		\item 
L'écart-type est: $\sigma(M_n) = \ds\sqrt{V(M_n)}$

$V(M_n)=\dfrac{V(Z)}{n}=\dfrac{2,07}{n}$ donc $\sigma(M_n) = \ds\sqrt{ \dfrac{2,07}{n}}$

On cherche $n$ tel que: $\ds\sqrt{ \dfrac{2,07}{n}}\leqslant 0,5$

$\ds\sqrt{ \dfrac{2,07}{n}}\leqslant 0,5
\iff \dfrac{2,07}{n} \leqslant 0,25
\iff \dfrac{2,07}{0,25}\leqslant n
\iff n\geqslant 8,28$

Donc l'écart type de $M_{n}$ est inférieur ou égal à 0,5 pour $n\geqslant 9$.

		\item Pour les valeurs trouvées en \textbf{b.}, on va montrer que la probabilité que\\
		 $6,3 \leqslant M_{n} \leqslant 8,3$ est supérieure ou égale à 0,75 .

\begin{list}{\textbullet}{}
\item $\begin{gathered}P\left (6,3 \leqslant M_{n} \leqslant 8,3\right )
= P\left (7,3 -1 \leqslant M_{n} \leqslant 7,3+1\right )
= P\left ( -1 \leqslant M_{n}-7,3 \leqslant 1\right )\\
\phantom{P\left (6,3 \leqslant M_{n} \leqslant 8,3\right )}
= P\left ( -1 \leqslant M_{n}-E(M_n) \leqslant 1\right )
= P\left (  \left |M_{n}-E(M_n) \strut \right | \leqslant 1\right )\\
\phantom{P\left (6,3 \leqslant M_{n} \leqslant 8,3\right )}
= 1- P\left (  \left |M_{n}-E(M_n)\strut \right | > 1\right )
\end{gathered}$

\item $P\left (  \left |M_{n}-E(M_n)\strut \right | > 1\right ) + P\left (  \left |M_{n}-E(M_n)\strut \right | = 1\right ) = P\left (  \left |M_{n}-E(M_n)\strut \right | \geqslant 1\right )$

donc
$P\left (  \left |M_{n}-E(M_n)\strut \right | > 1\right )  \leqslant P\left (  \left |M_{n}-E(M_n)\strut \right | \geqslant 1\right )$.

\item D'après l'inégalité de Bienaymé Tchebychev:\\
$P\left (  \left |X-\mu \strut\right | \geqslant \delta \right ) \leqslant \dfrac{V}{\delta^2}$
donc
$P\left (  \left |M_{n}-E(M_n) \strut\right | \geqslant 1\right ) \leqslant \dfrac{V(M_n)}{1^2}$.

On déduit donc:
$P\left (  \left |M_{n}-E(M_n) \strut\right | > 1\right ) \leqslant V(M_n)$.

\item D'après la question précédente, on peut prendre $n\geqslant 9$ donc\\
$V(M_n)=\dfrac{2,07}{n} \leqslant \dfrac{2,07}{9}$
et donc
$P\left (  \left |M_{n}-E(M_n) \strut\right | > 1\right ) \leqslant \dfrac{2,07}{9}$.

\item On déduit que:
$ 1- P\left (  \left |M_{n}-E(M_n)\strut \right | > 1\right ) \geqslant 1-\dfrac{2,07}{9}$.

$1-\dfrac{2,07}{9}=\dfrac{6,93}{9} \approx 0,77 \geqslant 0,75$
\end{list}		

On a donc démontré que, pour $n\geqslant 9$, 
$P\left (6,3 \leqslant M_{n} \leqslant 8,3\strut\right ) \geqslant0,75$.
	\end{enumerate}
\end{enumerate}
